\scanpage{144}

une base de $\mathscr{H}\otimes_{\mathbb{Z}}\mathbb{R}$ (\inserted{fibre par fibre} sur $\mathbb{R}$), la donnée d'une telle structure complexe équivaut à la donnée de
\[
\tau=z_2/z_1
\tag*{(3)}
\]
en tant que section du faisceau structural $\mathcal{O}_S$. (\uline{NB} \deleted{il faut,} pour une structure complexe donnée sur $\mathscr{H}\otimes_{\mathbb{Z}}\mathbb{R}$, $z_1$ et $z_2$ sont deux sections\inserted{-bases} du faisceau inversible obtenu, et $\tau$ est la section de $\underline{\mathcal{O}}_S$ qui fait passer de l'une à l'autre :
\[
z_2=\tau z_1\,.
\tag*{(3 bis)}
\]
% slanted
\margin{Noter que $\forall s\in S$, $\tau(s)\in\mathbb{C}\smallsetminus\mathbb{R}$}

Une structure \uline{de Torelli} sur une \deleted{\unsure{fam}} courbe elliptique relative $E$ sur $\unsure{S_a|S{,}}$ (correspondant à $\mathscr{H}$ sur $S$ $+$ structure complexe sur $\mathscr{H}\otimes_{\mathbb{Z}}\mathbb{R}$) est la donnée d'un isomorphisme
\[
t:\mathbb{Z}_S^2\xrightarrow{\sim}\mathscr{H}\,,
\tag*{(4)}
\]
\deleted{\unsure{ou bien|i.e.}} \inserted{\deleted{\unsure{donc de|deux d}} ou encore} de deux sections $z_1,z_2$ de $\mathscr{H}$, \deleted{\unsure{satisfais|distinguées}} formant une base de $\mathscr{H}$ en \unsure{tte} fibre. La catégorie des courbes elliptiques relatives sur $S$, avec structure de Torelli, est \deleted{\unsure{équivalente à}} \inserted{``rigide'' \deleted{et}} (pas d'autom. $\ne\mathrm{id}$ d'un objet) et l'ens des

% archive sheet number: 143
